%% ════════════════════════════════════════════════════════════════════════════
\section{System Architecture}
\label{sec:architecture}
%% ════════════════════════════════════════════════════════════════════════════

\subsection{Architecture Overview}
\label{subsec:arch_overview}

The system follows a \emph{hybrid modular monolith with event-driven
messaging} architecture.  The backend is a single FastAPI process; however,
it is internally structured around three decoupled event channels: an MQTT
ingestion pipeline, a WebSocket fan-out layer, and an on-demand AI
pipeline bus.  This avoids the operational overhead of microservices while
preserving the loose coupling that makes each subsystem independently
testable.

Figure~\ref{fig:arch_overview} shows the five-tier logical view: IoT
devices at the edge, a Mosquitto message broker, the FastAPI backend with
its service groups, the React frontend, and external integrations
(Moodle, Ollama).

\begin{figure}[H]
  \centering
  \includegraphics[width=\linewidth]{fig01_system_architecture}
  \caption{Full system layered architecture: IoT $\rightarrow$ Edge
    $\rightarrow$ Backend $\rightarrow$ Frontend $\rightarrow$ External.}
  \label{fig:arch_overview}
\end{figure}

Table~\ref{tab:tech_stack} summarises the technology choices for each
layer, and Table~\ref{tab:subsystems} lists each subsystem's primary
responsibility.

\begin{table}[H]
\centering
\caption{Technology Stack Summary}
\label{tab:tech_stack}
\begin{tabular}{@{}p{1.8cm}p{2.2cm}p{3.5cm}@{}}
\toprule
\textbf{Layer} & \textbf{Technology} & \textbf{Key Libraries} \\
\midrule
Firmware    & Arduino C++    & PubSubClient, DHT, LiquidCrystal\_I2C \\
Backend     & Python 3.11    & FastAPI, SQLAlchemy 2.0, aiomqtt, APScheduler, httpx \\
Database    & PostgreSQL 15  & Alembic (migrations) \\
Cache       & Redis 7        & redis-py async \\
AI          & Ollama phi3-mini & httpx AsyncClient \\
Face Rec.   & DeepFace       & OpenCV, FaceNet, numpy \\
Frontend    & React 18       & Vite, TailwindCSS, Recharts, axios \\
Infra       & Docker Compose & nginx, eclipse-mosquitto:2 \\
\bottomrule
\end{tabular}
\end{table}

\begin{table}[H]
\centering
\caption{Subsystem Responsibilities}
\label{tab:subsystems}
\begin{tabular}{@{}p{2.0cm}p{5.5cm}@{}}
\toprule
\textbf{Subsystem} & \textbf{Responsibility} \\
\midrule
IoT (ESP32)        & Sense environment; publish MQTT; actuate relays \\
Edge (RPi)         & Run all services locally; no cloud dependency \\
MQTT Broker        & Message bus between ESP32 and backend \\
Backend (FastAPI)  & Orchestration, REST + WebSocket APIs, business logic \\
Real-time Layer    & Fan-out of sensor/attendance/alert events to browsers \\
AI/Insights Layer  & At-risk detection, forecasting, SQL analytics \\
Frontend (React)   & Professor dashboard; demo mode when offline \\
Data Layer         & PostgreSQL (durable) + Redis (ephemeral live state) \\
External           & Moodle LMS sync; Ollama LLM inference \\
\bottomrule
\end{tabular}
\end{table}

%% ────────────────────────────────────────────────────────────────────────────
\subsection{IoT Layer — ESP32 Firmware}
\label{subsec:iot_layer}

\subsubsection{Sensor Acquisition}
The ESP32-WROOM-32 runs an Arduino sketch that reads four sensors on a
5-second cycle: the DHT21 temperature/humidity sensor (GPIO~4), the
MQ-135 air-quality ADC (GPIO~34, ADC1 channel~6), and the ACP014
digital sound sensor (GPIO~35).  GPIO~34 and~35 are input-only pins,
which avoids conflicts with the WiFi subsystem that disables ADC2.

\subsubsection{MQTT Publishing}
Each reading is serialised as a JSON payload
\texttt{\{value, unit, ts\}} and published to the corresponding MQTT
topic at QoS~0.  A heartbeat message (\texttt{status} topic, every
30~seconds) allows the backend to detect device offline events via a
TTL-based Redis presence key.

\subsubsection{Relay Actuation and Local Auto-Control}
The ESP32 subscribes to \texttt{classroom/room1/relay/ac} and
\texttt{classroom/room1/relay/lighting}.  On receipt, it drives the
corresponding GPIO output through a logic-level converter
(3.3\,V~$\leftrightarrow$\,5\,V) to the 4-channel opto-isolated relay
module.  An on-device auto-control flag (\texttt{acAutoMode}) implements
temperature thresholding independently of the backend---a deliberate
redundancy that maintains AC control if MQTT connectivity drops.

Figure~\ref{fig:esp32_wiring} shows the complete GPIO wiring.

\begin{figure}[H]
  \centering
  \includegraphics[width=\linewidth]{fig02_esp32_wiring}
  \caption{ESP32 IoT node: sensor inputs, logic-level converter, relay
    module, and LCD~16$\times$2 wiring.}
  \label{fig:esp32_wiring}
\end{figure}

%% ────────────────────────────────────────────────────────────────────────────
\subsection{Edge Layer — Raspberry Pi 4B}
\label{subsec:edge_layer}

The Raspberry Pi~4B (4~GB RAM) acts as the central edge hub, running
all Docker services on a single host.  In production, the RPi also
hosts the RPi Camera Module~v2, which feeds the face-recognition loop
at 2~fps via OpenCV \texttt{VideoCapture(0)}.

When \texttt{FACE\_RECOGNITION\_ENABLED=false} (Docker/dev mode), a
stub recognition loop fires synthetic attendance events using randomly
sampled enrolled students, exercising the full bidirectional
present$\leftrightarrow$absent evaluation logic without a physical
camera.

Figure~\ref{fig:rpi_topology} illustrates the runtime topology on the
Raspberry Pi~4B, including all Docker service ports and the WiFi link
to the ESP32 node.

\begin{figure}[H]
  \centering
  \includegraphics[width=\linewidth]{fig08_rpi_topology}
  \caption{RPi~4B production runtime topology: all services and their
    communication paths to the ESP32 classroom node and professor
    browser.}
  \label{fig:rpi_topology}
\end{figure}

%% ────────────────────────────────────────────────────────────────────────────
\subsection{Communication Layer — MQTT}
\label{subsec:mqtt_layer}

\subsubsection{Topic Architecture}
The Mosquitto~2 broker hosts a single hierarchical topic tree under the
prefix \texttt{classroom/\{room\_id\}/}.  Table~\ref{tab:mqtt_schema}
defines all topics, directions, and payload schemas.

\begin{table}[H]
\centering
\caption{MQTT Topic Schema}
\label{tab:mqtt_schema}
\begin{tabular}{@{}p{2.8cm}p{1.1cm}p{3.6cm}@{}}
\toprule
\textbf{Topic suffix} & \textbf{Dir.} & \textbf{Payload / Period} \\
\midrule
\texttt{sensors/temperature} & ESP$\rightarrow$RPi & \texttt{\{value,unit,ts\}} · 5~s \\
\texttt{sensors/humidity}    & ESP$\rightarrow$RPi & \texttt{\{value,unit,ts\}} · 5~s \\
\texttt{sensors/air\_quality}& ESP$\rightarrow$RPi & \texttt{\{value,unit,ts\}} · 5~s \\
\texttt{sensors/sound}       & ESP$\rightarrow$RPi & \texttt{\{value:0|1,ts\}} · 5~s \\
\texttt{status}              & ESP$\rightarrow$RPi & \texttt{\{online:true,ts\}} · 30~s \\
\texttt{relay/ac}            & RPi$\rightarrow$ESP & \texttt{\{action:on|off|auto\}} \\
\texttt{relay/lighting}      & RPi$\rightarrow$ESP & \texttt{\{action:on|off|auto\}} \\
\texttt{alerts}              & RPi$\rightarrow$ESP & \texttt{\{type,value\}} $\rightarrow$ LCD \\
\bottomrule
\end{tabular}
\end{table}

\subsubsection{Subscription Strategy}
The backend subscribes using the wildcard pattern
\texttt{classroom/+/sensors/\#} to receive sensor messages from any
room ID and any sensor sub-topic in a single subscription.  A separate
subscription on \texttt{classroom/+/status} handles heartbeat messages.
ESP32 nodes subscribe to exact topics for relay commands.

Figure~\ref{fig:mqtt_topics} shows the complete topic tree.

\begin{figure}[H]
  \centering
  \includegraphics[width=\linewidth]{fig03_mqtt_topic_tree}
  \caption{MQTT topic tree under \texttt{classroom/\{room\_id\}/}.
    Blue nodes are ESP32-published; green nodes are backend-published.}
  \label{fig:mqtt_topics}
\end{figure}

%% ────────────────────────────────────────────────────────────────────────────
\subsection{Backend — FastAPI Application}
\label{subsec:backend}

The backend is organised into seven logical service groups.
Figure~\ref{fig:backend_modules} shows their connections.

\begin{figure}[H]
  \centering
  \includegraphics[width=\linewidth]{fig04_backend_modules}
  \caption{Backend internal module groups and their data-flow
    connections.}
  \label{fig:backend_modules}
\end{figure}

\textbf{Group~1 — Ingestion \& Real-time} (\texttt{mqtt\_bridge.py}).
A persistent \texttt{aiomqtt.Client} subscriber loop connects to
Mosquitto with exponential backoff.  On each sensor message: (i)~Redis
SET with 5-minute TTL, (ii)~fire-and-forget PostgreSQL write,
(iii)~non-blocking put to \texttt{sensor\_event\_queue}.  Three
\texttt{asyncio.Queue} instances decouple producers from WebSocket
consumers.

\textbf{Group~2 — Session \& Attendance} (\texttt{api/sessions.py},
\texttt{recognition\_loop.py}).  Session lifecycle: \texttt{POST
/api/sessions/start} activates the recognition loop; \texttt{POST
/api/sessions/\{id\}/end} stops it, runs the final evaluation cycle,
and triggers Moodle sync.  Manual adjustments via
\texttt{PATCH /api/attendance/\{id\}} set \texttt{adjusted\_by} to
prevent overwrite by the recognition loop.

\textbf{Group~3 — Control} (\texttt{api/control.py}).  Relay commands
simultaneously write Redis (instant dashboard state) and publish MQTT
(device actuation).  Redis is the authoritative relay state; PostgreSQL
is not used for relay records.

\textbf{Group~4 — Alert Engine} (\texttt{services/alert\_engine.py}).
An \texttt{AsyncIOScheduler} job runs every 30~seconds, reads the Redis
sensor cache, evaluates thresholds (Table~\ref{tab:auto_control}), and
upserts \texttt{alerts} rows.  Deduplication prevents duplicate alerts
for the same unacknowledged condition.

\textbf{Group~5 — Insights \& Analytics.}  Pure SQL analytics: weekly
trend, day/hour heatmap, comfort score, AC effectiveness, temperature
versus attendance scatter, air-quality versus sound correlation.
Role-filtered: professors see only their courses.

\textbf{Group~6 — AI Pipelines} (Section~\ref{sec:ai}).

\textbf{Group~7 — Auth} (\texttt{api/auth.py}).  JWT (HS256), bcrypt
passwords, role-based access control; \texttt{REQUIRE\_AUTH=false} in
development.

\begin{table}[H]
\centering
\caption{Auto-Control Rules}
\label{tab:auto_control}
\begin{tabular}{@{}p{1.6cm}p{2.4cm}p{3.5cm}@{}}
\toprule
\textbf{Sensor} & \textbf{Condition} & \textbf{Action} \\
\midrule
Temperature & $> 28$\,°C + AC in auto & Turn AC ON \\
Temperature & $< 22$\,°C + AC in auto & Turn AC OFF \\
Air quality & $> 500$~ppm            & Alert (no relay) \\
Sound       & Silent $> 30$~min (active session) & Attendance anomaly alert \\
Device      & Heartbeat TTL expired  & Device-offline alert \\
\bottomrule
\end{tabular}
\end{table}

\begin{table}[H]
\centering
\caption{APScheduler Job Inventory}
\label{tab:scheduler_jobs}
\begin{tabular}{@{}p{2.8cm}p{1.1cm}p{3.6cm}@{}}
\toprule
\textbf{Job} & \textbf{Interval} & \textbf{Purpose} \\
\midrule
\texttt{check\_thresholds} & 30~s & Sensor threshold evaluation, relay auto-control, alert creation \\
\texttt{retry\_moodle\_sync} & 10~min & Re-send up to 10 queued session IDs to Moodle \\
\texttt{mock\_sensor\_publish} & 5~s & (\texttt{MOCK\_MODE}) Synthetic MQTT sensor messages \\
\texttt{mock\_heartbeat} & 30~s & (\texttt{MOCK\_MODE}) Synthetic status heartbeat \\
\bottomrule
\end{tabular}
\end{table}

Table~\ref{tab:api_endpoints} lists the primary REST endpoints.

\begin{table}[H]
\centering
\caption{Primary REST API Endpoints}
\label{tab:api_endpoints}
\begin{tabular}{@{}p{0.8cm}p{3.8cm}p{2.9cm}@{}}
\toprule
\textbf{Method} & \textbf{Path} & \textbf{Description} \\
\midrule
GET  & \texttt{/api/sensors/latest}            & Latest sensor values (Redis) \\
GET  & \texttt{/api/sensors/history}           & Time-range history (PG) \\
POST & \texttt{/api/sessions/start}            & Create + start session \\
POST & \texttt{/api/sessions/\{id\}/end}       & End session, sync Moodle \\
POST & \texttt{/api/control/ac}                & Set AC relay state \\
POST & \texttt{/api/control/lighting}          & Set lighting relay state \\
GET  & \texttt{/api/sessions/\{id\}/attendance}& Attendance roster \\
PATCH& \texttt{/api/attendance/\{id\}}         & Manual adjustment \\
POST & \texttt{/api/students/\{id\}/enroll-face}& Face enrollment \\
GET  & \texttt{/api/at-risk}                   & At-risk list + trigger pipeline \\
GET  & \texttt{/api/forecasting}               & Forecast list + trigger pipeline \\
WS   & \texttt{/ws/classroom/\{room\_id\}}     & Live event stream \\
\bottomrule
\end{tabular}
\end{table}

%% ────────────────────────────────────────────────────────────────────────────
\subsection{Real-Time Layer — WebSocket}
\label{subsec:websocket}

Figure~\ref{fig:ws_fanout} shows the asyncio fan-out mechanism.  Three
background drain tasks loop indefinitely, pulling events from their
respective queues and calling \texttt{connection\_manager.broadcast(room\_id,
event)}.  The \texttt{ConnectionManager} maintains a
\texttt{dict[room\_id $\rightarrow$ set[WebSocket]]} for room-scoped
delivery.

\begin{figure}[H]
  \centering
  \includegraphics[width=\linewidth]{fig05_websocket_fanout}
  \caption{Event fan-out: three asyncio queues drain independently into
    the WebSocket broadcast layer.}
  \label{fig:ws_fanout}
\end{figure}

On connect, the server immediately sends a \texttt{snapshot} message
containing current Redis sensor values, relay states, and device-online
status, ensuring the dashboard is never blank on page load.  A
30-second server-side keepalive (ping/pong) detects stale connections.
The frontend reconnects with exponential backoff (3~s initial, 30~s
cap).  Table~\ref{tab:ws_messages} defines all WebSocket message types.

\begin{table}[H]
\centering
\caption{WebSocket Message Types (Server $\rightarrow$ Client)}
\label{tab:ws_messages}
\begin{tabular}{@{}p{1.5cm}p{2.4cm}p{3.6cm}@{}}
\toprule
\textbf{Type} & \textbf{Trigger} & \textbf{Payload fields} \\
\midrule
\texttt{snapshot}   & On connect       & sensors, relay, device\_online \\
\texttt{sensor}     & MQTT ingestion   & room\_id, sensor\_type, value, unit \\
\texttt{attendance} & Recognition match & session\_id, student\_id, name, confidence, status \\
\texttt{alert}      & Threshold breach & alert\_type, room\_id, message, value \\
\texttt{ping}       & 30~s keepalive   & — \\
\bottomrule
\end{tabular}
\end{table}

%% ────────────────────────────────────────────────────────────────────────────
\subsection{Data Layer}
\label{subsec:data_layer}

\subsubsection{PostgreSQL Schema}
The database contains nine tables.
Figure~\ref{fig:er_diagram} shows the entity-relationship diagram.

\begin{figure}[H]
  \centering
  \includegraphics[width=\linewidth]{fig06_er_diagram}
  \caption{Entity-relationship diagram for all nine database tables.}
  \label{fig:er_diagram}
\end{figure}

Table~\ref{tab:db_schema} summarises each table's purpose and key
design decisions.

\begin{table}[H]
\centering
\caption{Database Schema Summary}
\label{tab:db_schema}
\begin{tabular}{@{}p{2.5cm}p{5.0cm}@{}}
\toprule
\textbf{Table} & \textbf{Key Design Points} \\
\midrule
\texttt{professors}        & UUID PK, bcrypt password, role ENUM (professor|admin) \\
\texttt{students}          & UUID PK, UNIQUE institutional student\_id \\
\texttt{courses}           & UUID PK, UNIQUE code, nullable professor FK \\
\texttt{sessions}          & display\_status computed (not stored) to reflect live state \\
\texttt{attendance\_records} & adjusted\_by guard: recognition loop never overwrites manual changes \\
\texttt{face\_encodings}   & 128-d float32 BYTEA; no raw images stored \\
\texttt{sensor\_readings}  & Time-series; queried for history and AI pipeline sensor averages \\
\texttt{at\_risk\_explanations} & UNIQUE student\_id; per\_course\_data as JSONB \\
\texttt{attendance\_forecasts} & UNIQUE course\_id; VARCHAR(30) for classification, not PG ENUM \\
\bottomrule
\end{tabular}
\end{table}

Key schema decisions: \texttt{display\_status} is computed at query time
because session liveness changes over time without a write.
\texttt{per\_course\_data} is JSONB because it is always read/written as
a unit with the parent row.  \texttt{trend\_classification} uses
\texttt{VARCHAR(30)} instead of a PostgreSQL ENUM to avoid
non-transactional \texttt{ALTER TYPE ADD VALUE} DDL on schema evolution.

\subsubsection{Redis Key Patterns}
Table~\ref{tab:redis_keys} defines all Redis key patterns, their
semantics, and TTL behaviour.

\begin{table}[H]
\centering
\caption{Redis Key Patterns}
\label{tab:redis_keys}
\begin{tabular}{@{}p{3.0cm}p{1.5cm}p{3.0cm}@{}}
\toprule
\textbf{Key pattern} & \textbf{TTL} & \textbf{Semantics} \\
\midrule
\texttt{classroom:\{r\}:sensors:\{t\}} & 300~s & Latest sensor value for type \texttt{t} \\
\texttt{classroom:\{r\}:online}        & 60~s  & Device presence (heartbeat-renewed) \\
\texttt{classroom:\{r\}:relay:\{d\}}   & None  & Relay state; persists across restarts \\
\texttt{at\_risk:pipeline:lock}        & 600~s & Prevents concurrent pipeline runs \\
\texttt{forecast:pipeline:lock}        & 1800~s & Prevents concurrent forecast runs \\
Moodle retry list                      & None  & Session IDs awaiting Moodle sync \\
\bottomrule
\end{tabular}
\end{table}

%% ────────────────────────────────────────────────────────────────────────────
\subsection{Frontend — React Dashboard}
\label{subsec:frontend}

The frontend is a single-page application (React~18, Vite, React
Router~v6).  Two context providers supply global state:
\texttt{AuthContext} (JWT token, professor profile) and
\texttt{SensorContext} (wrapping the \texttt{useLiveSensors} hook that
manages the WebSocket connection and demo-mode fallback).

The design system separates tokens from component classes:
\texttt{tokens.css} defines all CSS custom properties (colours,
spacing, radius, typography); \texttt{index.css} uses those variables
in component class rules.  TailwindCSS~v3 utility classes coexist but
cannot resolve CSS custom properties at build time, making the token
file split necessary.

Figures~\ref{fig:dashboard}--\ref{fig:control} show the three primary
professor-facing views.

\begin{figure}[H]
  \centering
  \includegraphics[width=\linewidth]{screenshot_dashboard}
  \caption{Dashboard page: live sensor sparklines, session control, and
    relay status cards.}
  \label{fig:dashboard}
\end{figure}

\begin{figure}[H]
  \centering
  \includegraphics[width=\linewidth]{screenshot_attendance}
  \caption{Attendance page: per-student status table with manual
    adjustment capability.}
  \label{fig:attendance}
\end{figure}

\begin{figure}[H]
  \centering
  \includegraphics[width=\linewidth]{screenshot_control}
  \caption{Control page: relay toggle controls with live sensor summary
    cards.}
  \label{fig:control}
\end{figure}

%% ────────────────────────────────────────────────────────────────────────────
\subsection{External Integrations}
\label{subsec:external}

\textbf{Moodle 4.x.}  The optional Moodle integration calls the Moodle
REST API (token-authenticated) on session end to write attendance
records.  Failures are enqueued in Redis; the alert engine retries up
to 10 items every 10~minutes.  Moodle runs as an optional Docker Compose
profile to avoid adding 600~MB of container weight to the default
development environment.

\textbf{Ollama / phi3-mini.}  The Ollama service exposes a REST API at
\texttt{http://ollama:11434}.  On backend startup, an async background
task checks for the phi3:mini model via \texttt{GET /api/tags} and pulls
it (\texttt{POST /api/pull}) if absent---a non-blocking operation that
takes approximately 2~GB of download on first run.  Both AI pipelines
share a single \texttt{httpx.AsyncClient} to reuse TCP connections and
avoid per-call TCP handshake overhead.

%% ────────────────────────────────────────────────────────────────────────────
\subsection{Deployment Architecture}
\label{subsec:deployment}

Figure~\ref{fig:docker_compose} shows the Docker Compose service
dependency graph.  PostgreSQL and Redis must pass health checks before
the backend starts.  The frontend starts after the backend is running,
so that the nginx proxy target is always reachable.

\begin{figure}[H]
  \centering
  \includegraphics[width=\linewidth]{fig07_docker_compose}
  \caption{Docker Compose service graph showing dependency order and
    container image sources.}
  \label{fig:docker_compose}
\end{figure}

Alembic migrations run automatically in the FastAPI lifespan context
(\texttt{upgrade head}) before any API routes are registered, ensuring
schema currency on every container restart without a manual migration
step.

%% ────────────────────────────────────────────────────────────────────────────
\subsection{Key Architectural Decisions}
\label{subsec:arch_decisions}

Table~\ref{tab:arch_decisions} captures the most significant
non-obvious architectural choices and their rationale.

\begin{table}[H]
\centering
\caption{Key Architectural Decisions}
\label{tab:arch_decisions}
\begin{tabular}{@{}p{2.5cm}p{5.0cm}@{}}
\toprule
\textbf{Decision} & \textbf{Rationale} \\
\midrule
FastAPI over Flask          & Native async for WebSocket + MQTT concurrent I/O \\
All services on RPi         & Avoids latency, cost, and internet dependency \\
asyncio-mqtt $\rightarrow$ aiomqtt & paho-mqtt~v2 broke the asyncio-mqtt API \\
DeepFace stub in Docker     & TensorFlow adds $\sim$600~MB; causes OOM on dev laptops \\
Redis for sensor state      & Sub-millisecond read on page load; no PostgreSQL hit \\
\texttt{display\_status} computed & Liveness is time-relative; storing it would require continuous writes \\
\texttt{tokens.css} split   & TailwindCSS~v3 cannot resolve CSS vars at build time \\
Glassmorphism removed       & GPU-expensive; caused rendering artifacts in Chrome \\
Inline SVG icons            & No icon-library dependency; only 10--12 icons needed \\
\texttt{per\_course\_data} JSONB & Always read/written as a unit; avoids join \\
Redis TTL not deleted on completion & Prevents re-run on every page refresh (natural cooldown) \\
1~LLM call/student          & N+1 calls: 14--16~min; batched: 2--3~min \\
Deterministic classification & LLM labels are unreliable for structured values; delta math is deterministic \\
\texttt{VARCHAR(30)} not PG ENUM & \texttt{ALTER TYPE ADD VALUE} is non-transactional in PostgreSQL \\
Marker rows ($<3$ sessions)  & Without a row, the frontend poll condition never resolves \\
No Redis cooldown in webcam endpoint & \texttt{last\_posted\_status} in the laptop script is the dedup layer; backend cooldown silently dropped transitions \\
\bottomrule
\end{tabular}
\end{table}
